\section{引言}

无人机集群在协同侦察、目标拦截、区域巡逻和空中对抗等任务中具有重要应用价值。与单机任务相比，集群任务更依赖多个平台之间的信息共享、角色分工和协同行为。在协同追逃场景中，多架追击无人机需要在动态目标机动和平台间安全约束下完成拦截；如果通信受限或协同策略不稳定，系统容易出现追击效率下降、无人机间碰撞或任务超时。

多智能体强化学习为无人机协同决策提供了端到端策略学习工具。集中训练、分散执行范式允许训练阶段使用全局信息，而执行阶段每个智能体仅依赖局部观测。MAPPO 等方法已经在多种合作式多智能体任务中表现出较强基线能力~\cite{yu2021mappo}。然而，普通 MAPPO 缺少显式关系建模机制，当通信半径受限时，不同无人机之间的相对空间关系和通信可达性难以被策略充分利用。图神经网络和图注意力机制为多智能体关系建模提供了自然表示~\cite{velickovic2017gat,malysheva2020magnet,liu2024gnn_marl}，但普通 GAT 主要从节点特征中计算注意力权重，并未显式利用无人机任务中关键的距离、方位、相对速度和通信约束。

现实无人机集群通常无法假设全局通信。通信半径、链路质量和带宽限制会改变智能体可获取的信息集合。已有通信学习和 GNN-based MARL 工作关注如何学习智能体间通信机制~\cite{singh2018ic3net,cuzin2026gnn_comm_survey}，UAV 通信与边缘计算场景中也已有图注意力和多智能体强化学习应用~\cite{feng2024gat_uav_comm,kim2024uav_mec_madrl}。但在无人机协同追逃任务中，通信约束并非单纯的邻接关系变化；距离、方位、相对速度和可达性会共同影响追击队形、安全间隔和目标包围过程。因此，如何将物理边语义嵌入多智能体策略，并在显式变化的通信半径下保持稳定协同，仍需要进一步研究。

本文从现实可实现角度出发，先在二维异构无人机协同追逃环境中验证有限通信图强化学习方法，而不直接进入完整 6DOF 空战、导弹、雷达和有人机协同系统。这样可以降低仿真复杂度，隔离算法贡献，并为后续迁移到 LAG/JSBSim 等平台提供可复用表示层。本文提出 EA-RG-MAPPO-S 方法，其核心由三部分组成：边特征增强角色图表示、基于 MAPPO 的集中训练分散执行优化、分阶段随机通信半径微调。

本文主要贡献如下：
\begin{enumerate}
    \item 提出一种边特征增强角色图策略表示，将追击无人机和目标建模为带角色语义的动态图节点，并将相对位置、距离、方位、相对速度和通信可达性引入图注意力得分，从而增强策略对任务相关物理关系的利用。
    \item 设计分阶段随机通信半径微调机制，使策略先学习固定通信半径下的基础协同能力，再适应变化通信拓扑，从而提升跨通信半径稳定性，而不是仅依赖单一通信半径训练。
    \item 构建有限通信异构无人机协同追逃实验，对 MAPPO、GAT-MAPPO 和多个消融版本进行多种子对比。300 回合每种子复评显示，所提方法在四个通信半径下将碰撞率控制在 0.054--0.086，并通过轨迹图、散点图和注意力热力图分析方法行为。
\end{enumerate}

需要说明的是，本文曾探索目标意图辅助分支，但诊断结果显示该分支当前 balanced accuracy 不足。因此，本文不将高精度目标意图识别作为主创新点，而将研究重点放在有限通信下的边特征角色图协同决策。
